\section{The scalar Tate calculation}
\label{sec:arithmetic-tate}

A local zeta integral separates into valuation shells, each weighted by
a power of the prime. The resulting geometric series contributes to a
single integral on the rational ideles. Its continuation comes from
Fourier summation on a rational lattice. The finite unit balls
select the integers from that lattice, and the real Gaussian supplies the
theta transformation. We construct the fields, measures, and lattice before
making this argument.

All vector spaces of test functions below have complex coefficients,
pointwise multiplication, and cohomological degree zero. At a finite place,
the locally constant functions on a fixed compact support with a fixed
open subgroup of constancy form a finite-dimensional space; their filtered
union has the locally convex inductive-limit topology. At the real place
we use the Schwartz topology, defined by the seminorms
\(\sup_x|x^a f^{(b)}(x)|\), \(a,b\geq0\).
The adelic test space below is the inductive limit of its finite restricted
tensor products with these topologies. These test algebras on the
noncompact fields and on the adeles have no pointwise multiplicative unit. The displayed integrals are scalar
linear functionals on their stated convergence domains; they are not
asserted to preserve pointwise products.
Integrals use positive Borel measures. A complex power of a positive real
number means \(a^s=\exp(s\log a)\), with the real logarithm.
We use countable product measures, Tonelli's theorem, dominated convergence,
and elementary real and complex analysis. The gamma function initially means
only Euler's integral
\[
 \Gamma(z)=\int_0^\infty e^{-u}u^{z-1}\,du,
 \qquad \operatorname{Re}z>0.
\]
No meromorphic continuation of \(\Gamma\) is needed in the proof.

\subsection{The local field and its unit ball}
\label{sec:adeles-local-field}

Fix a prime integer \(p\). Define the ring
\[
 \mathbb Z_p=\varprojlim_{n\geq1}\mathbb Z/p^n\mathbb Z
 =\{(a_n)_{n\geq1}:a_{n+1}\equiv a_n\pmod{p^n}\}.
\]
Addition and multiplication act on each coordinate. An integer maps to its
compatible residue classes. This map is injective, since an integer divisible
by every \(p^n\) is zero. Each element has unique digits
\[
 a=\sum_{j\geq0}a_jp^j,\qquad a_j\in\{0,\ldots,p-1\}.
\]
This notation specifies residues modulo \(p^n\). It does not assert
convergence in the real numbers.

Give each finite residue ring its discrete topology and \(\mathbb Z_p\) the induced
product topology. Its basic neighborhoods are the residue classes
\(a+p^n\mathbb Z_p\). Indeed, multiplication by \(p^n\) identifies \(\mathbb Z_p\)
with the elements whose first \(n\) digits vanish.
A diagonal subsequence argument shows compactness: successively fix each
digit, and then take the diagonal subsequence. The resulting subsequence
converges in every residue coordinate. This topology is metrizable by the
distance defined below, so sequential compactness implies compactness.
Each residue class is both open and closed.

An element \(u\in\mathbb Z_p\) is a unit precisely when its residue modulo \(p\)
is nonzero. In that case its inverse in each ring \(\mathbb Z/p^n\mathbb Z\)
exists by Bezout's identity. Uniqueness of the inverse makes these inverses
compatible, so they define \(u^{-1}\in\mathbb Z_p\).
For \(a\neq0\), its first nonzero digit therefore gives a unique expression
\[
 a=p^r u,\qquad r\in\mathbb N_0,\quad u\in\mathbb Z_p^\times.
\]
The product of two such expressions has exponent equal to the sum of their
exponents. In particular, \(\mathbb Z_p\) has no zero divisors.

Define \(\mathbb Q_p=\mathbb Z_p[1/p]\), the localization of this integral domain.
Every nonzero element now has a unique expression
\begin{equation}
 x=p^r u,\qquad r\in\mathbb Z,\quad u\in\mathbb Z_p^\times.
 \label{eq:adeles-valuation-decomposition}
\end{equation}
It is invertible. Thus \(\mathbb Q_p\) is a field and contains \(\mathbb Q\).
Put \(v_p(x)=r\), \(|x|_p=p^{-r}\), \(v_p(0)=+\infty\), and
\(|0|_p=0\). Divisibility gives
\[
 |xy|_p=|x|_p|y|_p,\qquad
 |x+y|_p\leq\max\{|x|_p,|y|_p\}.
\]
Consequently \(d_p(x,y)=|x-y|_p\) is a metric. On \(\mathbb Z_p\) it gives the
residue topology above. Addition and multiplication are continuous.
Inversion is continuous on \(\mathbb Q_p^\times\), since
\[
 |x^{-1}-y^{-1}|_p=\frac{|x-y|_p}{|x|_p|y|_p}
\]
and \(|y|_p=|x|_p\) whenever \(|y-x|_p<|x|_p\).

\begin{proposition}[Completion and local compactness]
\label{prop:adeles-completion}
The field \(\mathbb Q_p\) is complete for \(d_p\), and \(\mathbb Q\) is dense
in it. The sets \(a+p^r\mathbb Z_p\), with \(a\in\mathbb Q_p\) and \(r\in\mathbb Z\),
are compact open balls and form a neighborhood basis.
Thus \(\mathbb Q_p\) is a locally compact topological field.
\end{proposition}

\begin{proof}
A Cauchy sequence is bounded: its tail lies in a ball of radius one around
one of its terms, and its remaining terms form a finite set.
Multiplication by a sufficiently large power of \(p\) places the entire
sequence in \(\mathbb Z_p\). Each residue coordinate is eventually constant.
These constants are compatible and define its limit. This proves
completeness after scaling back. Truncated digit expansions approximate
every element of \(\mathbb Z_p\) by integers. Scaling proves density of
\(\mathbb Q\) in \(\mathbb Q_p\). The balls are translates and dilates of the
compact open set \(\mathbb Z_p\). The metric inequalities identify them with the
stated neighborhood basis.
\end{proof}

The notation \(\mathbb Q_p\) therefore agrees with the metric completion
of \(\mathbb Q\) for the usual rational \(p\)-adic absolute value.
No algebraic extension or choice of a root is involved.

\subsection{Measures, characters, and one Fourier calculation}
\label{sec:adeles-local-fourier}

On \(\mathbb Z_p\), give the digits independent uniform probability distributions.
The resulting measure assigns mass \(p^{-n}\) to every residue class modulo
\(p^n\). Translation and multiplication by a unit permute these residue
classes. They therefore preserve the measure, first on cylinder sets and
then on Borel sets by uniqueness of the product measure.
Also, multiplication by \(p\) scales this measure by \(p^{-1}\).

Extend the measure to \(\mathbb Q_p=\bigcup_{m\geq0}p^{-m}\mathbb Z_p\) by
\[
 \mu_p(E)=p^m\mu_p(p^mE),\qquad E\subset p^{-m}\mathbb Z_p.
\]
On the right, \(\mu_p\) denotes the probability measure on \(\mathbb Z_p\).
The scaling identity makes these prescriptions compatible as \(m\)
increases. They define a Borel measure on the union. Write it as \(dx_p\).
Translations preserve it, and
\begin{equation}
 \operatorname{vol}(a+p^r\mathbb Z_p)=p^{-r},\qquad
 d(bx)_p=|b|_p\,dx_p\quad(b\neq0).
 \label{eq:adeles-additive-scaling}
\end{equation}
The second identity follows from the unit and power-of-\(p\) cases.
This is an additive Haar measure: it is translation invariant, finite on
compact sets, and nonzero on open sets. Regularity follows from the
countable basis of compact open balls and the regular probability measure
on each \(p^{-m}\mathbb Z_p\). Its normalization determines it uniquely, since
the masses of all balls are then forced.

For \(x\in\mathbb Q_p\), its Laurent digit expansion has only finitely many
negative powers. Define
\[
 \{x\}_p=\sum_{r\leq j<0}a_jp^j
 \in\mathbb Z[1/p]\cap[0,1),\qquad
 \psi_p(x)=\exp(-2\pi i\{x\}_p).
\]
The fractional part is the unique element of this set congruent to \(x\)
modulo \(\mathbb Z_p\). To see uniqueness, note that
\(\mathbb Z[1/p]\cap\mathbb Z_p=\mathbb Z\). Two representatives in \([0,1)\)
cannot differ by a nonzero integer.
In particular, \(\{x+y\}_p-\{x\}_p-\{y\}_p\) is an integer.
Thus \(\psi_p:(\mathbb Q_p,+)\to\{z\in\mathbb C:|z|=1\}\) is a continuous
additive character, with kernel exactly \(\mathbb Z_p\).

For a locally constant, compactly supported function \(f:\mathbb Q_p\to\mathbb C\),
define its Fourier transform by
\[
 \widehat f(y)=\int_{\mathbb Q_p}f(x)\overline{\psi_p(xy)}\,dx_p.
\]
Such functions form the Schwartz--Bruhat space \(\mathcal S(\mathbb Q_p)\).

At \(p=2\) and \(y=1/2\), the even and odd residue classes each have
mass \(1/2\). The Fourier integral of \(\one_{\mathbb Z_2}\) is therefore
\[
 \frac12\cdot1+\frac12\cdot(-1)=0.
\]
At \(y\in\mathbb Z_2\), the same integral is one. The following identity
gives the full calculation, including every scale and translate.

\begin{proposition}[Fourier transform of a ball]
\label{prop:adeles-ball-fourier}
For \(a,y\in\mathbb Q_p\) and \(r\in\mathbb Z\),
\begin{equation}
 \widehat{\one_{a+p^r\mathbb Z_p}}(y)
 =p^{-r}\overline{\psi_p(ay)}\one_{p^{-r}\mathbb Z_p}(y).
 \label{eq:adeles-ball-fourier}
\end{equation}
In particular, \(\widehat{\one_{\mathbb Z_p}}=\one_{\mathbb Z_p}\).
Moreover, \(\widehat{\widehat f}(x)=f(-x)\) for every
\(f\in\mathcal S(\mathbb Q_p)\). Thus the additive measure of unit-ball mass one
is self-dual for this Fourier convention.
\end{proposition}

\begin{proof}
For \(y\in p^{-r}\mathbb Z_p\),
the character \(z\mapsto\psi_p(zy)\) is trivial on \(p^r\mathbb Z_p\).
Its integral over that ball is \(p^{-r}\).
If \(y\notin p^{-r}\mathbb Z_p\), then \(p^ry\notin\mathbb Z_p\).
Translation by \(p^r\) multiplies the integral by the nontrivial scalar
\(\overline{\psi_p(p^ry)}\). Translation invariance forces the integral
to vanish. Translation by \(a\) proves
\eqref{eq:adeles-ball-fourier}.

Every locally constant compactly supported function is a finite linear
combination of indicators of balls of one sufficiently small radius.
Indeed, finitely many neighborhoods of constancy cover its support.
Their finite refinement by balls of a common radius has this property.
Applying the Fourier transform twice to a ball indicator gives
\[
 p^{-r}\int_{p^{-r}\mathbb Z_p}
       \overline{\psi_p((a+x)y)}\,dy_p
 =\one_{a+p^r\mathbb Z_p}(-x).
\]
Linearity proves the assertion for the entire test space.
\end{proof}

Define a measure on \(\mathbb Q_p^\times\) by
\begin{equation}
 d^{\times} x_p=(1-p^{-1})^{-1}\frac{dx_p}{|x|_p}.
 \label{eq:adeles-multiplicative-measure}
\end{equation}
The scaling formula \eqref{eq:adeles-additive-scaling} proves multiplicative
invariance. The shell \(\mathbb Z_p^\times=\mathbb Z_p\setminus p\mathbb Z_p\) has additive
mass \(1-p^{-1}\). It therefore has multiplicative mass one.
Every shell \(p^r\mathbb Z_p^\times\) also has multiplicative mass one.
This is a locally finite regular Haar measure on the multiplicative group.

A quasi-character of a topological group is a continuous homomorphism to
\(\mathbb C^\times\). For every \(s\in\mathbb C\), the map
\(x\mapsto|x|_p^s\) is a quasi-character of \(\mathbb Q_p^\times\).

\begin{proposition}[The finite local integral]
\label{prop:adeles-finite-factor}
The integral
\[
 Z_p(s)=\int_{\mathbb Q_p^\times}\one_{\mathbb Z_p}(x)|x|_p^s\,d^{\times} x_p
\]
is absolutely convergent exactly when \(\operatorname{Re}s>0\). In that
half-plane,
\[
 Z_p(s)=\sum_{r\geq0}p^{-rs}=(1-p^{-s})^{-1}.
\]
With additive measure instead,
\[
 \int_{\mathbb Z_p}|x|_p^{s-1}\,dx_p
 =\frac{1-p^{-1}}{1-p^{-s}},\qquad\operatorname{Re}s>0.
\]
The value assigned to the additive integrand at zero has no effect.
\end{proposition}

\begin{proof}
The disjoint valuation decomposition
\eqref{eq:adeles-valuation-decomposition} reduces the first integral to the
geometric series. Its absolute integral is \(\sum_{r\geq0}p^{-r\operatorname{Re}s}\).
The additive shell \(p^r\mathbb Z_p^\times\) has mass
\((1-p^{-1})p^{-r}\), which gives the second formula.
The singleton \(\{0\}\) has measure zero since it lies in balls of mass
\(p^{-n}\) for every \(n\).
\end{proof}

These measure normalizations are the rational unramified case of Tate's
conventions \cite{Tate1950}*{\S2.2, pp.~309--311, and \S2.3, pp.~312--313}.
Tate's additive character is the complex conjugate of the one fixed here.
Conjugating the character leaves these self-dual measures unchanged.
The computations above establish the formulas directly, including the prime two.

\subsection{The real Gaussian and Poisson summation}
\label{sec:adeles-real-poisson}

At the real place, put
\[
 \psi_\infty(x)=e^{2\pi ix},\qquad dx_\infty=dx,
 \qquad d^{\times} x_\infty=\frac{dx}{|x|}.
\]
Here \(dx\) is Lebesgue measure, and \(d^{\times} x_\infty\) is a measure on
\(\mathbb R^\times\). Both signs occur. The real Schwartz space
\(\mathcal S(\mathbb R)\) consists of smooth functions whose derivatives
decay faster than every inverse power. Its Fourier transform is
\[
 \widehat f(y)=\int_{\mathbb R}f(x)e^{-2\pi ixy}\,dx.
\]
Differentiation under the integral and integration by parts show that this
map preserves \(\mathcal S(\mathbb R)\).

\begin{lemma}[Gaussian Fourier transform and Mellin integral]
\label{lem:adeles-gaussian}
For \(t>0\), define \(g_t(x)=e^{-\pi t^2x^2}\). Then
\[
 \widehat g_t(y)=t^{-1}e^{-\pi y^2/t^2}.
\]
For \(\operatorname{Re}s>0\),
\[
 Z_\infty(s)=\int_{\mathbb R^\times}e^{-\pi x^2}|x|^s\,d^{\times} x_\infty
 =\pi^{-s/2}\Gamma(s/2)=:\GammaR(s).
\]
The latter integral is absolutely convergent exactly in this half-plane.
\end{lemma}

\begin{proof}
The positive Gaussian integral has value one. Indeed, its square is the
integral of \(e^{-\pi(x^2+y^2)}\) over the plane. Polar coordinates give
\(2\pi\int_0^\infty e^{-\pi r^2}r\,dr=1\).
Set \(G(y)=\widehat g_1(y)\). Differentiation and integration by parts give
\(G'(y)=-2\pi yG(y)\), with \(G(0)=1\).
The boundary term vanishes by Gaussian decay. Thus
\(G(y)=e^{-\pi y^2}\). Substitution \(x\mapsto x/t\) gives the scaled
identity.

Evenness and the substitution \(u=\pi x^2\) give
\[
 2\int_0^\infty e^{-\pi x^2}x^{s-1}\,dx
 =\pi^{-s/2}\int_0^\infty e^{-u}u^{s/2-1}\,du.
\]
Gaussian decay controls infinity. At zero, absolute convergence is
equivalent to \(\operatorname{Re}s>0\).
\end{proof}

\begin{proposition}[The local factors]
\label{prop:adeles-local-factors}
\label{prop:local-tate}
With the specified additive character and multiplicative measures, for
\(\operatorname{Re}s>0\),
\[
 Z_p(\one_{\mathbb Z_p},s)=(1-p^{-s})^{-1},\qquad
 Z_\infty(e^{-\pi x^2},s)=\GammaR(s)=\pi^{-s/2}\Gamma(s/2).
\]
Here \(Z_v(f,s)=\int_{\mathbb Q_v^\times}f(x)|x|_v^s\,d^{\times} x_v\),
where \(\mathbb Q_\infty=\mathbb R\).
\end{proposition}

\begin{proof}
These are Proposition~\ref{prop:adeles-finite-factor} and
Lemma~\ref{lem:adeles-gaussian} with their test functions displayed.
\end{proof}

For completeness, the Gaussian also proves real Fourier inversion with
this measure. If \(f\in\mathcal S(\mathbb R)\) and \(\epsilon>0\), Fubini
and Lemma~\ref{lem:adeles-gaussian} give
\[
 \int_{\mathbb R}\widehat f(y)e^{2\pi ixy}e^{-\pi\epsilon y^2}\,dy
 =\int_{\mathbb R}f(v)\epsilon^{-1/2}
       e^{-\pi(x-v)^2/\epsilon}\,dv.
\]
The right side tends to \(f(x)\), since the positive kernels have mass one
and their mass outside every fixed neighborhood of zero tends to zero.
The left side tends to the unregularized integral by dominated convergence,
because \(\widehat f\) is integrable. Hence
\(\widehat{\widehat f}(x)=f(-x)\), which fixes the self-dual normalization.

\begin{proposition}[Real Poisson summation]
\label{prop:adeles-real-poisson}
For \(f\in\mathcal S(\mathbb R)\), \(d>0\), and \(r\in\mathbb R\),
\begin{equation}
 \sum_{n\in\mathbb Z}f(r+dn)
 =\frac1d\sum_{m\in\mathbb Z}
       \widehat f(m/d)e^{2\pi irm/d}.
 \label{eq:adeles-real-poisson}
\end{equation}
Both series converge absolutely. For fixed \(d\), they converge uniformly
as \(r\) ranges over a compact set.
\end{proposition}

\begin{proof}
Periodize \(f\) by \(P(x)=\sum_n f(x+dn)\). Rapid decay of all derivatives
gives uniform convergence on \([0,d]\), including each differentiated
series. Thus \(P\) is smooth and \(d\)-periodic.
Its \(m\)th Fourier coefficient is
\[
 \frac1d\int_0^dP(x)e^{-2\pi imx/d}\,dx
 =\frac1d\int_{\mathbb R}f(x)e^{-2\pi imx/d}\,dx
 =\frac1d\widehat f(m/d).
\]
Absolute integrability justifies the interchange and the partition into
intervals. The coefficients decay faster than every inverse power, so
their Fourier series converges absolutely and uniformly.

To verify equality with \(P\), one may use the elementary Fejer kernels
\[
 K_N(u)=\frac1N\left|\sum_{j=0}^{N-1}e^{2\pi iju/d}\right|^2.
\]
They are nonnegative and satisfy \(d^{-1}\int_0^dK_N(u)\,du=1\).
Outside a fixed neighborhood of \(d\mathbb Z\), the geometric-series
formula bounds them by a constant times \(N^{-1}\).
Uniform continuity therefore makes convolution with \(K_N\), using
\(du/d\), converge uniformly to any continuous periodic function.
A function with zero Fourier coefficients has all these convolutions zero.
Apply this observation to the difference between \(P\) and its uniformly
convergent Fourier series, then evaluate at \(r\).
\end{proof}

\subsection{The adele ring and the rational fundamental domain}
\label{sec:adeles-fundamental-domain}

Define the finite adele ring and the full adele ring by
\[
 \mathbb A_f=\prod_p'(\mathbb Q_p,\mathbb Z_p)
 =\{(x_p)_p:x_p\in\mathbb Z_p\text{ for all but finitely many }p\},
 \qquad \Adele=\mathbb R\times\mathbb A_f.
\]
Ring operations are componentwise. The prime on the product imposes the
displayed integrality condition. Put \(\widehat{\mathbb Z}=\prod_p\mathbb Z_p\).
The restricted product topology has basic open sets
\[
 \prod_{p\in S}U_p\times\prod_{p\notin S}\mathbb Z_p,
\]
where \(S\) is finite and \(U_p\subset\mathbb Q_p\) is open.
These sets cover \(\mathbb A_f\), and their translates form the same basis.
In particular, \(\widehat{\mathbb Z}\) is compact and open. Compactness follows from
the countable product of the compact metrizable spaces \(\mathbb Z_p\), or from
the same diagonal subsequence argument used for their digits.
Addition and multiplication are continuous: at the finitely many
exceptional primes this follows from local continuity, and outside them
\(\mathbb Z_p\) is closed under both operations. Thus \(\Adele\) is a locally
compact Hausdorff topological ring.

Every rational number belongs to \(\mathbb Z_p\) at all primes except those
dividing its denominator. Hence the diagonal embedding
\(\mathbb Q\hookrightarrow\Adele\) is defined.

Give \(\widehat{\mathbb Z}\) its product probability measure. On each open subring
\[
 A_S=\mathbb R\times\prod_{p\in S}\mathbb Q_p\times\prod_{p\notin S}\mathbb Z_p
\]
use the product of Lebesgue measure, the finitely many local additive
measures, and the remaining product probability measure.
These measures agree on overlaps because \(\operatorname{vol}(\mathbb Z_p)=1\).
An increasing exhaustion by finite sets of primes defines a measure
\(dx=dx_\infty\prod_pdx_p\) on \(\Adele\).
Translation by an adele preserves it: each relevant set and translation
lies in one sufficiently large \(A_S\), where product invariance applies.
It is regular and locally finite, as follows on each open \(A_S\) from
the corresponding product measures. This constructs the additive Haar
measure without an unspecified constant.

\begin{proposition}[The rational lattice]
\label{prop:adeles-rational-lattice}
The diagonal subgroup \(\mathbb Q\subset(\Adele,+)\) is discrete and closed.
The set
\begin{equation}
 \mathcal F=[0,1)\times\widehat{\mathbb Z}
 \label{eq:adeles-fundamental-domain}
\end{equation}
contains exactly one representative of each coset in \(\Adele/\mathbb Q\).
It has additive measure one, and its closure is compact.
For every nonnegative Borel function \(H\),
\[
 \int_{\Adele}H(x)\,dx
 =\int_{\mathcal F}\sum_{q\in\mathbb Q}H(x+q)\,dx.
\]
The same formula holds for an absolutely integrable complex function.
\end{proposition}

\begin{proof}
For \(x\in\Adele\), the sum \(r=\sum_p\{x_p\}_p\) is a finite rational
sum. At a fixed prime \(p\), each term with another prime in its denominator
belongs to \(\mathbb Z_p\). Consequently \(x_p-r\in\mathbb Z_p\) for every \(p\).
Choose the integer \(n=\lfloor x_\infty-r\rfloor\). Then
\(x-(r+n)\in\mathcal F\).

A rational number integral at every prime is an integer: in a reduced
fraction, any prime dividing the denominator would violate integrality.
Two representatives in \(\mathcal F\) that differ by a rational number
therefore differ by an integer. Their real coordinates force that integer
to be zero. This proves uniqueness.

The neighborhood \((-1,1)\times\widehat{\mathbb Z}\) meets \(\mathbb Q\) only at zero,
so the subgroup is discrete. It is also closed. Indeed, choose a smaller
neighborhood \(V\) with \(V-V\subset(-1,1)\times\widehat{\mathbb Z}\).
Each translate of \(V\) contains at most one rational point. A point in
the closure of \(\mathbb Q\) must then equal that point, by the Hausdorff
property. The product measure gives \(\operatorname{vol}(\mathcal F)=1\),
and \([0,1]\times\widehat{\mathbb Z}\) is compact. Finally, the countably many sets
\(\mathcal F+q\) form a disjoint Borel partition. Translation invariance
and Tonelli's theorem prove the integral formula.
\end{proof}

Define
\begin{equation}
 \psi(x)=\psi_\infty(x_\infty)\prod_p\psi_p(x_p)
 =\exp\!\left(2\pi i\left(x_\infty-\sum_p\{x_p\}_p\right)\right).
 \label{eq:adeles-character}
\end{equation}
Only finitely many factors differ from one. On each \(A_S\), the character
is a finite product of continuous characters. It is therefore continuous
on \(\Adele\). For \(q\in\mathbb Q\), the rational number
\(q-\sum_p\{q\}_p\) is integral at every prime, hence belongs to
\(\mathbb Z\). Thus \(\psi(q)=1\), and \(\psi\) descends to
\(\Adele/\mathbb Q\).

The annihilator of the diagonal \(\mathbb Q\) under the pairing
\((x,y)\mapsto\psi(xy)\) is exactly the diagonal \(\mathbb Q\).
Here is a direct verification. Subtract a rational number from an element
\(y\) of the annihilator to put it in \(\mathcal F\).
Testing against integers first gives \(y_\infty=0\), since
\(\psi(ny)=e^{2\pi iny_\infty}\). Testing against \(p^{-k}\), for each
prime \(p\) and \(k\geq1\), then gives \(y_p\in p^k\mathbb Z_p\).
All other finite components contribute one. Hence every \(y_p\) is zero.
Conversely, \(\psi(qr)=1\) for rational \(q,r\).
The pairing therefore uses the rational lattice with its precise dual and
covolume. These are the rational cases of Tate's additive global data
\cite{Tate1950}*{Theorems~4.1.3--4.1.4, pp.~330--331}.

\subsection{Fourier summation on the adeles}
\label{sec:adeles-poisson}

Define \(\mathcal S(\Adele)\) to be the finite linear span of functions
\[
 F(x)=f_\infty(x_\infty)\prod_p f_p(x_p),
 \qquad f_\infty\in\mathcal S(\mathbb R),\quad
 f_p\in\mathcal S(\mathbb Q_p),
\]
where \(f_p=\one_{\mathbb Z_p}\) for all but finitely many primes.
This is the Schwartz--Bruhat test space used here.
Every such function is integrable for the additive Haar measure.
Fubini and the local ball calculation give the product Fourier transform
\[
 \widehat F(y)=\int_{\Adele}F(x)\overline{\psi(xy)}\,dx.
\]
For a fixed \(y\), only finitely many local integrals differ from the
unit-ball integral. The latter vanishes when its frequency is nonintegral.
The transform preserves the displayed test space, and its square is
\(F(x)\mapsto F(-x)\), by the proved local inversions.

\begin{lemma}[Elementary finite-adele test functions]
\label{lem:adeles-elementary-tests}
Every \(F\in\mathcal S(\Adele)\) is a finite sum of functions
\[
 f(x_\infty)\one_{r+d\widehat{\mathbb Z}}(x_f),
 \qquad f\in\mathcal S(\mathbb R),\quad r\in\mathbb Q,\quad
 d\in\mathbb Q_{>0}.
\]
\end{lemma}

\begin{proof}
Decompose each of the finitely many exceptional local factors into ball
indicators, using the proof of Proposition~\ref{prop:adeles-ball-fourier}.
Each resulting finite factor is a constant times an indicator of a coset
\(a_f+\prod_p p^{r_p}\mathbb Z_p\), with only finitely many \(r_p\neq0\).
Set \(d=\prod_p p^{r_p}>0\). Multiplication by the other prime factors
acts by units on \(\mathbb Z_p\), so this subgroup is \(d\widehat{\mathbb Z}\).
The proof of Proposition~\ref{prop:adeles-rational-lattice}, applied to
\(a_f/d\), supplies \(b\in\mathbb Q\) with \(a_f/d-b\in\widehat{\mathbb Z}\).
Then \(r=db\) is the required rational representative.
\end{proof}

For a rational \(q\neq0\), prime factorization proves the product formula
\begin{equation}
 |q|_\infty\prod_p|q|_p=1.
 \label{eq:adeles-product-formula}
\end{equation}
In particular, the additive finite measure of \(d\widehat{\mathbb Z}\) is \(d^{-1}\)
for \(d\in\mathbb Q_{>0}\).

\begin{theorem}[Adelic Poisson summation over the rationals]
\label{thm:adeles-poisson}
For every \(F\in\mathcal S(\Adele)\),
\begin{equation}
 \sum_{q\in\mathbb Q}F(q)=\sum_{q\in\mathbb Q}\widehat F(q).
 \label{eq:adeles-poisson}
\end{equation}
Both series converge absolutely. The measure and character are exactly
\eqref{eq:adeles-additive-scaling} and \eqref{eq:adeles-character}.
\end{theorem}

\begin{proof}
By Lemma~\ref{lem:adeles-elementary-tests}, it suffices to take
\(F(x)=f(x_\infty)\one_{r+d\widehat{\mathbb Z}}(x_f)\).
A rational \(q\) satisfies \(q-r\in d\widehat{\mathbb Z}\) precisely when
\((q-r)/d\in\mathbb Z\). Thus
\[
 \sum_{q\in\mathbb Q}F(q)=\sum_{n\in\mathbb Z}f(r+dn).
\]
Put \(\psi_f(x_f)=\prod_p\psi_p(x_p)\).
The local Fourier formula and the product formula give
\[
 \widehat F(y)
 =d^{-1}\widehat f(y_\infty)
       \overline{\psi_f(ry_f)}\one_{d^{-1}\widehat{\mathbb Z}}(y_f).
\]
For rational \(q\), the finite support condition is \(q=m/d\), with
\(m\in\mathbb Z\). Since \(rq\) is rational, \(\psi(rq)=1\), whence
\(\overline{\psi_f(rq)}=e^{2\pi irq}\). Therefore
\[
 \sum_{q\in\mathbb Q}\widehat F(q)
 =d^{-1}\sum_{m\in\mathbb Z}\widehat f(m/d)e^{2\pi irm/d}.
\]
Proposition~\ref{prop:adeles-real-poisson} identifies the two sides and
proves absolute convergence. A finite sum completes the proof.
\end{proof}

In particular, put
\[
 \phi(x)=e^{-\pi x_\infty^2}\one_{\widehat{\mathbb Z}}(x_f).
\]
The local formulas prove \(\widehat\phi=\phi\).
For \(t>0\), let \(a_t=(t,(1)_p)\) be the adele with real coordinate
\(t\) and all finite coordinates one. Applying
Theorem~\ref{thm:adeles-poisson} to \(x\mapsto\phi(a_tx)\) gives
\begin{equation}
 \Theta(t)=t^{-1}\Theta(t^{-1}),\qquad
 \Theta(t):=\sum_{n\in\mathbb Z}e^{-\pi n^2t^2}.
 \label{eq:adeles-theta}
\end{equation}
Here the finite unit balls impose \(q\in\mathbb Z\), and the Fourier
transform is \(t^{-1}\phi(a_{t^{-1}}x)\).
This proves the needed summation identity with its scaling constant.

\subsection{Ideles and absolute convergence of the Euler integral}
\label{sec:adeles-ideles}

The idele group is
\[
 \Adele^\times=\mathbb R^\times\times\prod_p'(\mathbb Q_p^\times,\mathbb Z_p^\times).
\]
It is the group of units of the adele ring as a set: both \(x\) and
\(x^{-1}\) are adeles precisely when the finite coordinates are units
at almost every prime. Its topology is the multiplicative restricted
product topology. In particular,
\(U=\prod_p\mathbb Z_p^\times\) is a compact open subgroup of the finite ideles.
This topology matters. It is not being inferred from the subspace topology
of the adele ring.

On each open product with finitely many exceptional primes, take the
product of the local multiplicative Haar measures. The remaining unit
groups carry their probability measures. These products agree on overlaps
because each \(\mathbb Z_p^\times\) has mass one. They define
\[
 d^{\times} x=\frac{dx_\infty}{|x_\infty|}\prod_pd^{\times} x_p.
\]
The same exhaustion argument used for additive measure proves local
finiteness, regularity, and multiplicative invariance.
For an idele define
\[
 |x|_{\Adele}=|x_\infty|\prod_p|x_p|_p\in\mathbb R_{>0}.
\]
This finite product is a continuous homomorphism. It is trivial on
\(\mathbb Q^\times\) by \eqref{eq:adeles-product-formula}.
Consequently \(x\mapsto|x|_{\Adele}^s\) is a quasi-character of the idele
group, trivial on \(\mathbb Q^\times\), for every \(s\in\mathbb C\).

Consider the scalar integral
\begin{equation}
 Z(\phi,s)=\int_{\Adele^\times}\phi(x)|x|_{\Adele}^s\,d^{\times} x.
 \label{eq:adeles-global-integral}
\end{equation}
It is an integral over the idele group, with \(\phi\) restricted from
the adele ring.

\begin{proposition}[Absolute Euler product]
\label{prop:adeles-euler-product}
The integral \eqref{eq:adeles-global-integral} is absolutely convergent
exactly when \(\operatorname{Re}s>1\). In this half-plane,
\begin{equation}
 Z(\phi,s)=\GammaR(s)\prod_p(1-p^{-s})^{-1}
 =\GammaR(s)\sum_{n\geq1}n^{-s}.
 \label{eq:adeles-euler}
\end{equation}
The product and series converge locally uniformly there.
\end{proposition}

\begin{proof}
Let \(\sigma=\operatorname{Re}s\). On the support of \(\phi\) in the
ideles, the finite valuations are nonnegative and only finitely many are
nonzero. For a finite set \(S\) of primes, restrict the valuations outside
\(S\) to zero. For \(\sigma>0\), Tonelli and the finite local calculation
give the absolute integral on this set as
\[
 \GammaR(\sigma)\prod_{p\in S}(1-p^{-\sigma})^{-1}
 =\GammaR(\sigma)
   \sum_{\substack{n\geq1\\p\mid n\Rightarrow p\in S}}n^{-\sigma}.
\]
The last identity follows by multiplying finitely many absolutely
convergent geometric series and using unique factorization.
An increasing exhaustion of the primes exhausts the support.
Every positive integer eventually occurs. Monotone convergence therefore
identifies the full absolute integral with
\(\GammaR(\sigma)\sum_{n\geq1}n^{-\sigma}\), allowing the value infinity.
The series converges exactly for \(\sigma>1\), by comparison with
\(\int_1^\infty u^{-\sigma}\,du\).
For \(\sigma\leq0\), already the real integral on the finite-unit subset
diverges at zero.

In the convergent half-plane, dominated convergence justifies the same
exhaustion for the complex integrand. Finite factorization then proves
\eqref{eq:adeles-euler}. On a compact subset with
\(\operatorname{Re}s\geq1+\delta\), the Dirichlet series has majorant
\(\sum_n n^{-1-\delta}\). The finite Euler products converge uniformly
by the same majorant and exhaustion. Equivalently, their logarithms have
the absolutely convergent expansion
\[
 \sum_p\sum_{k\geq1}\frac{p^{-ks}}{k},\qquad
 \sum_p\sum_{k\geq1}\frac{|p^{-ks}|}{k}
 \leq\frac{1}{1-2^{-1-\delta}}\sum_{n\geq2}n^{-1-\delta}.
\]
This also verifies the usual absolute convergence of the Euler product.
\end{proof}

\begin{proposition}[An explicit idele class domain]
\label{prop:adeles-idele-domain}
Every idele has a unique expression
\begin{equation}
 x=q\,(t,u),\qquad q\in\mathbb Q^\times,\quad t\in\mathbb R_{>0},
 \quad u\in U.
 \label{eq:adeles-idele-domain}
\end{equation}
Here \((t,u)\) denotes an idele, and \(t=|x|_{\Adele}\).
In these coordinates the measure is counting measure on
\(\mathbb Q^\times\), multiplied by \(dt/t\) and probability measure
\(du\) on \(U\). Thus
\[
 \mathbb Q^\times\backslash\Adele^\times
 \simeq\mathbb R_{>0}\times U
\]
as topological groups, with quotient measure \((dt/t)du\).
\end{proposition}

\begin{proof}
Let \(r_p=v_p(x_p)\), with only finitely many nonzero values, and set
\[
 q=\operatorname{sgn}(x_\infty)\prod_p p^{r_p}.
\]
Then \(u_p=x_p/q\) is a unit at every prime, while
\(t=x_\infty/q>0\). The product formula gives
\(t=|x|_{\Adele}\). Conversely, the finite valuations and real sign recover
\(q\) from any expression of the stated form, so it is unique.
The subset \(D=\mathbb R_{>0}\times U\) is an open subgroup of the
idele group, and the sets \(qD\) are its disjoint open cosets.
Restriction of \(d^{\times} x\) to \(D\) is exactly \((dt/t)du\).
Multiplicative invariance proves the measure statement on each coset.
The quotient map restricted to \(D\) is a continuous open group
isomorphism. This proves the topological assertion as well.
\end{proof}

In particular, the subgroup of idele classes of norm one has volume one.
There is no extra factor from the two real signs: the rational factor
\(q\) in \eqref{eq:adeles-idele-domain} already records that sign.

\subsection{Continuation from the Gaussian theta identity}
\label{sec:adeles-continuation}

For \(\operatorname{Re}s>1\), the preceding domain and absolute convergence
allow unfolding over \(\mathbb Q^\times\):
\begin{align}
 Z(\phi,s)
 &=\int_U\int_0^\infty
      \sum_{q\in\mathbb Q^\times}\phi(q(t,u))t^s\,\frac{dt}{t}\,du
 \notag\\
 &=\int_0^\infty(\Theta(t)-1)t^s\,\frac{dt}{t}.
 \label{eq:adeles-mellin-theta}
\end{align}
Indeed, since every \(u_p\) is a unit, the finite factors require
\(q\in\mathbb Z\setminus\{0\}\). Both positive and negative integers
occur. The sum is independent of \(u\), whose total measure is one.
This proves the Mellin identity directly from the idele integral.

For \(t\geq1\), the elementary inequality \(n^2\geq n\) gives
\begin{equation}
 0<\Theta(t)-1
 =2\sum_{n\geq1}e^{-\pi n^2t^2}
 \leq\frac{2e^{-\pi t^2}}{1-e^{-\pi t^2}}
 \leq\frac{2}{1-e^{-\pi}}e^{-\pi t^2}.
 \label{eq:adeles-theta-bound}
\end{equation}
It follows that
\[
 E(s)=\int_1^\infty(\Theta(t)-1)t^s\,\frac{dt}{t}
\]
is entire. More explicitly, on every compact set of \(s\)-values and for
every integer \(j\geq0\), the differentiated integrand is bounded by
a constant times \(e^{-\pi t^2}t^{M-1}(\log t)^j\), for some fixed
\(M\). This bound is integrable on \([1,\infty)\). Differentiation under
the integral therefore proves holomorphy and all derivative formulas.

\begin{theorem}[The completed rational zeta integral]
\label{thm:adeles-tate}
\label{thm:tate}
With the character, measures, and test function defined above, the
absolutely convergent integral on \(\operatorname{Re}s>1\) equals
\[
 Z(\phi,s)=\GammaR(s)\zeta(s)=:\Lambda(s),\qquad
 \zeta(s)=\sum_{n\geq1}n^{-s}.
\]
Its unique meromorphic continuation to \(\mathbb C\) is
\begin{equation}
 \Lambda(s)=\frac1{s-1}-\frac1s+E(s)+E(1-s).
 \label{eq:adeles-continuation}
\end{equation}
It has precisely two poles, both simple, with
\[
 \operatorname*{Res}_{s=1}\Lambda(s)=1,\qquad
 \operatorname*{Res}_{s=0}\Lambda(s)=-1,
 \qquad \Lambda(s)=\Lambda(1-s).
\]
Consequently,
\[
 \xi_{\mathbb R}(s)=\frac12s(s-1)\Lambda(s)
\]
is entire and satisfies \(\xi_{\mathbb R}(s)=\xi_{\mathbb R}(1-s)\).
Its values at both \(0\) and \(1\) are \(1/2\).
\end{theorem}

\begin{proof}
Proposition~\ref{prop:adeles-euler-product} proves the identity in the
initial half-plane. Split \eqref{eq:adeles-mellin-theta} at \(t=1\).
For \(0<t<1\), the proved theta identity gives
\[
 \Theta(t)-1=t^{-1}(\Theta(t^{-1})-1)+t^{-1}-1.
\]
When \(\operatorname{Re}s>1\), each of the following integrals converges
absolutely. Substitution \(v=t^{-1}\) gives
\begin{align*}
 \int_0^1(\Theta(t)-1)t^s\,\frac{dt}{t}
 &=\int_1^\infty(\Theta(v)-1)v^{1-s}\,\frac{dv}{v}
   +\int_0^1(t^{s-1}-t^s)\,\frac{dt}{t}\\
 &=E(1-s)+\frac1{s-1}-\frac1s.
\end{align*}
This proves \eqref{eq:adeles-continuation} on the original domain.
Since \(E\) is entire, its right side defines a meromorphic function on
\(\mathbb C\). The identity theorem gives uniqueness of the continuation.
The rational part has residues \(1\) and \(-1\), and the remaining terms
are entire. Thus neither pole cancels, and no other pole occurs.

Replacing \(s\) by \(1-s\) exchanges the two entire terms and preserves
the rational part. This proves the functional equation. Finally,
\[
 \xi_{\mathbb R}(s)
 =\frac12+\frac12s(s-1)\bigl(E(s)+E(1-s)\bigr)
\]
proves the asserted entireness, symmetry, and endpoint values.
\end{proof}

The transformation in \eqref{eq:adeles-continuation} is the theta-integral
argument in Riemann's paper \cite{Riemann1859}*{p.~3 of the transcription}.
Tate's global theorem extends this mechanism to number fields and Hecke
characters \cite{Tate1950}*{Main Theorem~4.4.1, pp.~339--341}.
Here every summation and domain formula
needed for the rational Gaussian integral has been proved explicitly.
The continuation is a meromorphic scalar function. An operator family or
cohomological determinant requires additional constructions on its own
carrier.
